\documentclass[final,5p,times,twocolumn]{elsarticle}

\usepackage{amssymb}
\usepackage{amsmath}
\usepackage{booktabs}
\usepackage{graphicx}
\usepackage{array}
\usepackage{multirow}
\usepackage{hyperref}
\usepackage{microtype}

\input{generated/results_macros.tex}

\journal{Physiological Measurement}

\begin{document}

\begin{frontmatter}

\title{How far can a subject-independent sEMG array be reduced? A prespecified non-inferiority test of cross-fitted attribution for electrode selection}

\author[inst1]{Minh-Tri Ta}
\author[inst2,inst3]{Hoang-Khang Phan}
\author[inst4]{Le-Giang Tran\corref{cor1}}
\cortext[cor1]{Corresponding author. Email: tranlegiang@iuh.edu.vn}

\affiliation[inst1]{organization={College of Engineering and Computer Science, VinUniversity},
            city={Hanoi},
            country={Vietnam}}
\affiliation[inst2]{organization={Department of Biomedical Engineering, Faculty of Applied Science, Ho Chi Minh City University of Technology (HCMUT)},
            addressline={268 Ly Thuong Kiet Street, District 10},
            city={Ho Chi Minh City},
            country={Vietnam}}
\affiliation[inst3]{organization={Vietnam National University Ho Chi Minh City},
            addressline={Linh Trung Ward, Thu Duc City},
            city={Ho Chi Minh City},
            country={Vietnam}}
\affiliation[inst4]{organization={Faculty of Electronics Technology, Industrial University of Ho Chi Minh City},
            city={Ho Chi Minh City},
            country={Vietnam}}

\begin{abstract}
\textit{Objective.} To test, under decision rules fixed in advance, whether explainability-guided electrode selection can reduce a forearm sEMG array without meaningful loss of subject-independent gesture recognition.
\textit{Approach.} We used 43 participants of GRABMyo (16 monopolar forearm electrodes, 16 active gestures) and leave-one-subject-out evaluation, in which the held-out participant never influenced scaling, training or ranking. Electrodes were ranked by TreeSHAP attributions cross-fitted over participant-disjoint inner folds. A stability-regularised variant (SR-CFA) penalised attribution that varied across those folds. A frozen analysis plan fixed the participant as the statistical unit, a 0.02 Macro-F1 non-inferiority margin, 10,000-resample cluster bootstrap intervals and Holm correction.
\textit{Main results.} With all electrodes, XGBoost reached a Macro-F1 of \FullXGB\ (95\% CI \FullXGBLCB--\FullXGBUCB). The smallest non-inferior subset was $k^* = \KStar$ electrodes ($\bar\Delta = \DeltaKTwelve$, 95\% CI $\DeltaKTwelveLCB$ to $\DeltaKTwelveUCB$). One fewer electrode was not non-inferior, and $k^*$ was \KStarMarginPtOne\ and \KStarMarginPtThree\ under the 0.01 and 0.03 sensitivity margins. Attribution-ranked subsets outperformed random subsets at 4, 6, 8 and 12 electrodes (largest gain $\HTwoKSixDiff$ at six). Deletion with retraining confirmed that the ranking was faithful. The stability penalty changed almost no rankings and gave no measurable benefit. In exploratory comparisons, a data-free evenly spread layout matched the attribution ranking at most sizes.
\textit{Significance.} A population-level array of this design tolerates removal of about a quarter of its monopolar electrodes within a stated engineering tolerance, not more. Cross-fitted attribution identifies which positions matter, but it offers little beyond even spatial coverage for choosing sparse layouts.
\end{abstract}

\begin{keyword}
Surface electromyography \sep gesture recognition \sep electrode selection \sep SHAP \sep cross-fitting \sep non-inferiority \sep leave-one-subject-out
\end{keyword}

\end{frontmatter}

\section{Introduction}
\label{sec:introduction}

Surface electromyography (sEMG) pattern recognition is the dominant approach to multifunction myoelectric control~\cite{hudgins1993new,englehart2003robust,scheme2011electromyogram,farina2014extraction}. Recognition generally improves with more recording sites, but every electrode adds cost, wiring, contact area and a point of failure in a wearable device~\cite{vujaklija2016new}. This makes it attractive to identify a reduced set of electrodes that preserves recognition performance~\cite{phinyomark2012feature,liu2014emg,altimemy2013classification,stango2015spatial,qu2021reduce}.

Three methodological problems limit what published channel-reduction results can tell a device designer. First, most reductions are selected and evaluated within the same users. A layout chosen for one person need not transfer to the next, and subject-independent performance is systematically lower than user-dependent performance~\cite{phinyomark2018emg,cote2019deep}. If data from a test user influences scaling, ranking or tuning, the reported subset is optimistically biased. Second, feature-attribution methods such as Shapley additive explanations~\cite{lundberg2017unified,lundberg2020local} are increasingly used to rank channels~\cite{salmi2025shapaad}, but attributions can be unstable across data splits~\cite{visani2022statistical,kelodjou2023shaping,xiang2026roshap}, and a ranking computed on the evaluation data is itself a form of leakage. Third, the chosen number of channels is usually read off an accuracy curve by eye, with no pre-stated tolerance and no statement of what ``no meaningful loss'' means~\cite{schuirmann1987comparison,walker2011understanding}. A margin chosen after the curve is seen cannot support a non-inferiority claim~\cite{campbell2018noninferiority}.

This study tests explainability-guided channel reduction under conditions designed to remove these three sources of optimism. Each held-out participant is excluded from every step that shapes the model or the channel ranking. Channel rankings are computed by cross-fitting: TreeSHAP attributions are always evaluated on participants the attributing model did not see. Every decision rule was fixed in a statistical analysis plan written before any sparse-channel result existed. This includes the non-inferiority margin ($\delta = 0.02$ Macro-F1), the participant as the statistical unit, the bootstrap, the multiplicity correction, and the definition of the smallest acceptable channel count $k^*$. We also test a stability-regularised variant of the ranking (SR-CFA), which penalises channels whose attribution varies across disjoint groups of training participants.

The study is primarily a test, not a method proposal. We asked four prespecified questions on the 43-participant GRABMyo dataset~\cite{pradhan2022multiday}:
(H1) whether some $k < 16$ is non-inferior to all 16 channels;
(H2) whether attribution-ranked subsets beat random subsets of the same size;
(H3) whether the stability penalty improves on attribution magnitude alone;
(H4) whether the attribution is faithful, i.e.\ whether removing highly ranked channels hurts more than removing random or low-ranked ones.
We report every answer, including the negative ones. Exploratory analyses covered cross-session transfer, a bipolar re-derivation of the montage, and a methodological replication on NinaPro DB5~\cite{pizzolato2017comparison}. These are reported as exploratory.

\section{Related Work}
\label{sec:related_work}

\paragraph{Channel reduction for myoelectric control}
Channel and feature reduction for sEMG pattern recognition has been approached with feature-ranking and wrapper methods~\cite{phinyomark2012feature,liu2014emg}, with channel selection on multi-channel and high-density arrays~\cite{altimemy2013classification,stango2015spatial}, and with sparse-representation methods~\cite{qu2021reduce}. Large-scale reviews of the field emphasise that most evaluations are user-dependent. They also note that population-level (cross-user) models behave differently from per-user models~\cite{phinyomark2018emg}. Deep cross-user models and transfer learning are one response to this gap~\cite{cote2019deep}. Our study sits at the other end of the spectrum: a conventional feature-based classifier, a strictly subject-independent protocol, and a question about hardware rather than model capacity.

\paragraph{Attribution-guided selection and attribution stability}
TreeSHAP computes Shapley-value attributions for tree ensembles~\cite{lundberg2017unified,lundberg2020local}. Attribution has been used to reduce channels in other biosignal tasks, for example EEG auditory attention decoding~\cite{salmi2025shapaad}. Several works show that attributions can change across resampling and neighbourhood choices~\cite{visani2022statistical,kelodjou2023shaping}. A recent preprint treats Shapley attributions as distributions across resamples and proposes a robustness-adjusted importance score~\cite{xiang2026roshap}. That idea is close to the stability penalty tested here, and we do not claim the penalty itself as new. The general idea of preferring features that are selected consistently across subsamples goes back to stability selection~\cite{meinshausen2010stability}. What this study adds is the evaluation design. Attribution is cross-fitted over participants inside each outer training fold. The stability penalty is tested as a prespecified hypothesis against magnitude alone, and the selected subsets are judged by a prespecified non-inferiority rule.

\paragraph{Faithfulness of attributions}
A ranking that looks plausible need not reflect what the model uses. Deletion-based tests remove the features an explanation marks as important and measure the effect~\cite{samek2017evaluating}. Retraining after removal (ROAR) avoids confounding the explanation with distribution shift of the unretrained model~\cite{hooker2019benchmark}. We use the retraining variant with channels as the unit.

\paragraph{Electrode shift and multi-day use}
Electrode displacement and day-to-day variation degrade myoelectric pattern recognition~\cite{hargrove2006effect,kaufmann2010fluctuating}, and adaptation methods have been proposed to counter the resulting covariate shift~\cite{vidovic2016improving,muceli2014extracting}. GRABMyo was recorded over three days~\cite{pradhan2022multiday}, which allows an exploratory cross-session analysis. However, the evaluation here is subject-independent throughout. Its cross-session contrasts therefore combine a session effect with the cross-user gap, and they are not a pure test of donning and doffing.

\section{Methods}
\label{sec:methods}

\begin{figure*}[t]
\centering
\includegraphics[width=\textwidth]{figures/fig1_workflow.pdf}
\caption{Study design. For every held-out participant $p$, all preprocessing statistics, model fitting and channel ranking use only the other participants $R_{-p}$. Rankings are cross-fitted: $R_{-p}$ is split into five participant-disjoint inner folds, and attributions are computed only on inner-held-out participants. Subsets of $k$ channels are then retrained from scratch and tested on $p$. Decision rules (margin, statistical unit, bootstrap, multiplicity, $k^*$) were frozen before any sparse-channel result existed.}
\label{fig:workflow}
\end{figure*}

\subsection{Datasets}

\paragraph{GRABMyo (primary)}
GRABMyo v1.1.0~\cite{pradhan2022multiday} contains recordings from \GRABMyoParticipants\ able-bodied participants (\GRABMyoMale\ male, \GRABMyoFemale\ female; age \GRABMyoAgeMean\ $\pm$ \GRABMyoAgeSD\ years) on \GRABMyoSessions\ different days. Signals were recorded with an EMGUSB2+ amplifier (OT Bioelettronica) at 2048~Hz. In each session participants performed 16 hand and wrist gestures plus rest, with 7 trials of 5~s each. We analysed the 16 forearm electrodes. These form two circumferential rings of 8 electrodes (F1--F8 and F9--F16), 2~cm apart along the forearm, with electrode $j$ of one ring aligned with electrode $j+8$ of the other. Mean circumference at the ring position was \GRABMyoCircMean\ $\pm$ \GRABMyoCircSD~cm, giving about \GRABMyoArcSpacing~cm between neighbouring electrodes in a ring. The recordings are monopolar, so the selection unit in the primary analysis is a single electrode. All \GRABMyoFilesVerified\ distributed files matched the dataset's published SHA-256 checksums. All \GRABMyoRecords\ records were processed, with \GRABMyoRecordsFailed\ failures. The primary task is the 16 active gestures (rest excluded), for which chance Macro-F1 is about 0.063.

\paragraph{NinaPro DB5 (exploratory methodological replication)}
DB5~\cite{pizzolato2017comparison} provides \DBFiveSubjects\ able-bodied subjects (\DBFiveMale\ male, \DBFiveFemale\ female; age \DBFiveAgeMean\ $\pm$ \DBFiveAgeSD\ years, from the files' metadata) recorded with two Myo armbands (16 channels, 200~Hz). We used exercise~1 (12 finger movements, 6 repetitions each), segmented with the refined \texttt{restimulus}/\texttt{rerepetition} labels. The DB5 values are 8-bit integer codes, not physical units. The spectra show no powerline peak, and their power extends to the Nyquist frequency. DB5 was analysed separately and never pooled with GRABMyo. Its channel numbers bear no anatomical correspondence to GRABMyo's, so only the behaviour of the procedure is compared across datasets.

\subsection{Preprocessing and features}
GRABMyo signals were band-pass filtered (20--450~Hz, 4th-order zero-phase Butterworth) and notch filtered at 50~Hz ($Q = 30$). A spectral audit chose 50~Hz over 60~Hz: mean band-power ratio was 10.35 at 50~Hz and 0.62 at 60~Hz. Windows were 250~ms with 50\% overlap and never crossed trial boundaries. From each channel we extracted ten features: mean absolute value (MAV), root mean square, variance, waveform length, zero crossings (ZC), slope sign changes (SSC), integrated EMG (IEMG), Willison amplitude (WAMP), mean frequency and median frequency. This gives a 160-dimensional vector per window, \GRABMyoWindows\ windows in total and \GRABMyoActiveWindows\ from active gestures. The ZC, SSC and WAMP amplitude threshold was set to $10^{-2}$~mV, about 1.15 times the median absolute first difference of the signal. This single label-free scalar was computed from the pooled recordings, not per fold. That is a minor departure from strict fold-wise estimation, which we note as a limitation. An earlier threshold 87 times smaller had saturated WAMP and was discarded. For DB5 we used a 10--95~Hz band-pass, no notch, 250~ms windows (50 samples) and a threshold of 7 ADC units, derived in the same way.

\subsection{Classifiers and leave-one-subject-out evaluation}
All evaluation is leave-one-subject-out (LOSO): 43 folds for GRABMyo and 10 for DB5. In fold $p$, z-scoring statistics, model fitting and channel ranking use only $R_{-p}$. The test participant contributes only to the score. XGBoost used 300 trees, depth 6, learning rate 0.1, row and column subsampling of 0.8, and the \texttt{hist} method. It was trained on all \FoldTrainingWindows\ outer-training windows. For compute reasons, the RBF support vector machine ($C = 1$, $\gamma$ = \texttt{scale}) used a stratified 40,000-window subsample and the random forest (200 trees, minimum leaf size 2) a 200,000-window subsample. On one calibration fold, halving the SVM subsample from 40,000 to 20,000 windows lowered Macro-F1 from \SVMCalibFortyK\ to \SVMCalibTwentyK. The SVM had therefore not saturated, and its result is a lower bound. No hyperparameter was tuned on any outer test participant. The primary metric is participant-level Macro-F1 over the 16 active gestures; accuracy and balanced accuracy were also recorded. XGBoost was prespecified as the primary model for the channel analyses on grounds of compute.

\subsection{Cross-fitted channel attribution and the SR-CFA score}
Within each outer fold, $R_{-p}$ was split into $K = 5$ participant-disjoint inner folds. For inner fold $j$, an XGBoost model was trained on the other inner folds (capped at a 200,000-window stratified subsample and z-scored with inner-training statistics only). Exact TreeSHAP attributions~\cite{lundberg2020local} were then computed for up to 20,000 windows of the inner-held-out participants. Per-window channel importance averaged absolute attributions over the 16 classes and the 10 features of the channel,
\begin{equation}
\Phi_{i,c} = \frac{1}{G F}\sum_{g=1}^{G}\sum_{f=1}^{F} \lvert \phi_{i,g,c,f} \rvert ,
\end{equation}
and was averaged over windows and normalised to sum to one across channels, giving $I_{j,c}$. Across the five inner folds, $A_c = \operatorname{median}_j I_{j,c}$ measures attribution magnitude. $U_c = \operatorname{median}_j \lvert I_{j,c} - A_c \rvert$ measures its instability across disjoint groups of training participants. Channels were ranked by
\begin{equation}
Q_c(\lambda) = A_c - \lambda\, U_c , \qquad \lambda \in \{0, 0.25, 0.5, 1\}.
\label{eq:srcfa}
\end{equation}
$\lambda = 0$ is cross-fitted SHAP magnitude; $\lambda > 0$ is the stability-regularised score (SR-CFA). A leakage assertion verified in every fold that no inner-held-out participant appeared in its inner training set. Unit tests verified that a participant-specific signal planted in the data could not raise that participant's score.

The frozen plan required $\lambda$ to be chosen on outer-training data only but did not specify a criterion, and no selection procedure was implemented. We therefore select no $\lambda$ and report all four (a post-hoc specification, Section~\ref{sec:deviations}).

\subsection{Sparse sensor curve and non-inferiority (H1)}
For every outer fold, $\lambda$ and $k \in \{1,\ldots,16\}$, the top-$k$ channels of the fold's own ranking were used to retrain XGBoost from scratch. Each model used the fold's fixed 200,000-window stratified outer-training subsample and was tested on the held-out participant. The $k = 16$ model of the same curve is the reference, so every paired difference $\Delta_{p,k} = \mathrm{F1}_{p,k} - \mathrm{F1}_{p,16}$ compares models with identical training budgets. Non-inferiority at margin $\delta$ requires the lower bound of the 95\% participant-cluster bootstrap interval of $\bar\Delta_k$ to exceed $-\delta$. We define $k^* = \min\{k < 16 : \mathrm{LCB}_{95}(\bar\Delta_k) > -\delta\}$, recorded as ``none'' if no $k$ qualifies. The primary margin $\delta = 0.02$ Macro-F1 is an engineering tolerance, not a clinically validated threshold. $\delta = 0.01$ and $0.03$ are prespecified sensitivity margins.

\subsection{Comparison with other selection methods (H2)}
At $k \in \{4, 6, 8, 12\}$ we compared six rankings, each computed on outer-training data only:
\begin{itemize}
\item SR-CFA ($\lambda = 0.5$);
\item cross-fitted SHAP ($\lambda = 0$);
\item permutation importance, from a model fitted on a 20,000-window subsample and evaluated on 8,000 further windows with 3 repeats;
\item mutual information, estimated on the same subsample;
\item a geometric baseline, which selects electrodes by farthest-point sampling on the $2 \times 8$ montage;
\item random subsets: 100 draws per participant at each $k$, and 500 at the prespecified $k^*$.
\end{itemize}
The value $\lambda = 0.5$ for this comparison had no recorded rationale. Because the P4 comparison includes $\lambda = 0$ as well, both rankings are reported. Drawing hundreds of random subsets per participant was only tractable with a lighter comparison model (XGBoost with 50 trees, trained on a 150,000-window stratified subsample), applied identically to every method. Its absolute scores are therefore lower than those of the sparse curve and are compared only within the comparison. H2 contrasts each attribution ranking with the participant's mean over its random subsets (Holm over 8 tests). We also report where the SR-CFA subset falls within each participant's random distribution.

\subsection{Stability regularisation (H3) and faithfulness (H4)}
H3 compared each $\lambda > 0$ with $\lambda = 0$ on the sparse curve at $k \in \{4, 6, 8, 12\}$ (Holm over 12 tests). For H4 we used deletion with retraining~\cite{hooker2019benchmark}. For $n \in \{2, 4, \ldots, 14\}$ retained channels, the model was retrained after removing either the $16 - n$ highest-ranked channels (keeping the $n$ lowest) or the $16 - n$ lowest-ranked (keeping the top $n$). Random removal is the mean over the random subsets above, which exist for $n \in \{4, 6, 8, 12\}$. All three arms used the same comparison model and training subsample. As a check on protocol equivalence, the top-$n$ arm selects the same electrodes as the H2 cross-fitted SHAP subsets in all \FaithProtocolCells\ matching participant-by-$k$ cells. Its Macro-F1 agrees with them to within a mean absolute difference of \FaithProtocolMeanDiff\ (maximum \FaithProtocolMaxDiff). The residual differences are numerical: z-scoring was applied per subset rather than once over all features. A faithful ranking gives top-removed $<$ random $<$ bottom-removed. Paired contrasts were Holm-corrected within H4.

\subsection{Statistics}
The participant is the statistical unit ($n = 43$; $n = 10$ for DB5). Windows are never treated as independent. Confidence intervals are 95\% percentile intervals from 10,000 cluster-bootstrap resamples of participants. Paired comparisons use the Wilcoxon signed-rank test with the matched-pairs rank-biserial correlation $r_{rb}$ as effect size. Holm correction is applied within each confirmatory family. Exploratory analyses are uncorrected and labelled as such. No participant was excluded.

\subsection{Exploratory analyses}
\paragraph{Ranking stability} We measured the agreement of rankings computed from single participants' attributions, from the five inner folds, and from the 43 outer folds. Outer folds share about 98\% of their training participants, so their agreement is only an upper bound.

\paragraph{Cross-session} XGBoost (full configuration) was trained on session~1 of $R_{-p}$ (CS1) and tested on each session of $p$. It was also trained on sessions 1 and 2 (CS2) and tested on session 3 of $p$. CS2 has roughly twice the training data of CS1, so the two conditions differ in both session diversity and quantity.

\paragraph{Bipolar montage} To express the selection unit as a conventional bipolar sensor, we derived $B_j = F_j - F_{j+8}$ ($j = 1, \ldots, 8$) and re-extracted features. The count-feature threshold was re-derived for this signal ($5 \times 10^{-3}$~mV), because the bipolar signal is about half the monopolar amplitude. We then repeated the LOSO baseline, the cross-fitted ranking and the sparse curve, with the $k = 8$ model as reference.

\paragraph{DB5} The LOSO baseline, cross-fitted ranking ($K = 5$) and sparse curve were repeated on DB5, using all outer-training windows and the $k = 16$ model as reference.

\paragraph{Fixed layout} In a post-hoc analysis, a single 12-electrode layout was derived from a cross-fitted ranking ($\lambda = 0.5$) on the 22 odd-numbered participants. Models with 16 and with 12 electrodes were trained on those participants and tested on the 21 even-numbered participants. This is one split, not a LOSO estimate.

\paragraph{Feature redundancy} IEMG equals 512 times MAV for these windows. We repeated the comparison-model LOSO with IEMG removed and compared both Macro-F1 and channel attribution.

\subsection{Deviations from the analysis plan}
\label{sec:deviations}
The following changes were made after results existed. None alters a prespecified decision rule.
\begin{enumerate}
\item $\lambda$ was not selected; see above.
\item The H2--H4 families were specified in detail (tests and $k$ sets) after the sparse curve and selection-baseline results had been computed.
\item The comparison model for H2 and H4 was reduced from 150 to 50 trees for compute reasons before any baseline result existed.
\item An earlier faithfulness run mixed arms trained with different budgets and was discarded, and all arms were recomputed under one protocol.
\item An earlier bipolar run used a count threshold inherited from the monopolar signal, and was discarded and recomputed.
\item The fixed-layout analysis was not planned.
\item The window-length sensitivity and the 17-class task, both listed as exploratory in the plan, were not run.
\end{enumerate}

\section{Results}
\label{sec:results}

\subsection{Full-channel subject-independent performance}
With all 16 electrodes, XGBoost reached a Macro-F1 of \FullXGB\ (95\% CI \FullXGBLCB--\FullXGBUCB) over the 43 held-out participants. The RBF-SVM reached \FullSVM\ (\FullSVMLCB--\FullSVMUCB) and the random forest \FullRF\ (\FullRFLCB--\FullRFUCB) (Table~\ref{tab:full_channel}, Fig.~\ref{fig:full}). XGBoost and the SVM did not differ: the paired difference was $\XGBvsSVMDiff$ (95\% CI $\XGBvsSVMLCB$ to $\XGBvsSVMUCB$), $p_{\mathrm{Holm}} \XGBvsSVMP$, $r_{rb} = \XGBvsSVMR$. Both exceeded the random forest ($p_{\mathrm{Holm}} \XGBvsRFP$ and $\SVMvsRFP$; $r_{rb} = \XGBvsRFR$ and $\SVMvsRFR$). The SVM result is a lower bound because of its training subsample (Methods). The spread between participants was far larger than any difference between models: XGBoost ranged from \FullXGBMin\ to \FullXGBMax\ across participants.

\begin{table*}[t]
\centering
\small
\caption{Full 16-channel leave-one-subject-out performance on GRABMyo (43 participants, 16 active gestures; chance Macro-F1 $\approx$ 0.063). CI: participant cluster bootstrap.}
\label{tab:full_channel}
\begin{tabular}{llcccc}
\toprule
Model & Training windows & Macro-F1 & 95\% CI & SD & Range \\
\midrule
XGBoost & all ($\approx$550k) & 0.542 & 0.517--0.567 & 0.085 & 0.385--0.702 \\
RBF-SVM & 40k & 0.538 & 0.513--0.562 & 0.085 & 0.383--0.718 \\
Random forest & 200k & 0.502 & 0.475--0.528 & 0.090 & 0.346--0.670 \\
\bottomrule
\end{tabular}
\par\smallskip\parbox{0.95\linewidth}{\footnotesize Paired comparisons (Wilcoxon signed-rank, Holm over 3): Random forest $-$ RBF-SVM = $-$0.0362 (95\% CI $-$0.0444 to $-$0.0284), $p_{\mathrm{Holm}}$ $<$ 0.001, $r_{rb}$ = -0.97; Random forest $-$ XGBoost = $-$0.0408 (95\% CI $-$0.0475 to $-$0.0340), $p_{\mathrm{Holm}}$ $<$ 0.001, $r_{rb}$ = -0.98; RBF-SVM $-$ XGBoost = $-$0.0045 (95\% CI $-$0.0107 to +0.0015), $p_{\mathrm{Holm}}$ = 0.154, $r_{rb}$ = -0.25.}
\end{table*}


\begin{figure*}[t]
\centering
\includegraphics[width=\textwidth]{figures/fig3_full_channel.pdf}
\caption{Full-channel LOSO performance (43 participants; each point is one held-out participant; grey lines join the same participant; black bar = median; diamond = mean with 95\% cluster-bootstrap CI). (a) Three classifiers. (b) Paired XGBoost $-$ SVM difference per participant, sorted; grey band = 95\% CI of the mean difference.}
\label{fig:full}
\end{figure*}

\subsection{H1: a 12-electrode subset is non-inferior under the prespecified margin}
Channel attribution was strongly concentrated and highly reproducible across outer folds. The same four electrodes (\TopFourSet) formed the top four in \TopFourFolds\ of 43 folds, and they carried \TopFourShare\% of the summed attribution. The eight lowest-ranked electrodes carried \BottomEightShare\% (Fig.~\ref{fig:attribution}). Electrodes \DroppedFourSet\ were the four excluded at $k = 12$ in \DroppedFourFolds\ of 43 folds.

Performance declined smoothly as electrodes were removed (Fig.~\ref{fig:curve}, Table~\ref{tab:sparse_ni}). The $k = 16$ reference of the curve (\CurveKSixteen) differs from the full-data result above because every curve model uses the same fixed 200,000-window training subsample. At $\delta = 0.02$, the smallest non-inferior subset was $k^* = \KStar$ electrodes: $\bar\Delta_{12} = \DeltaKTwelve$ (95\% CI $\DeltaKTwelveLCB$ to $\DeltaKTwelveUCB$). At $k = 11$ the lower bound ($\DeltaKElevenLCB$) crossed the margin. Every $k$ from 12 to 15 was non-inferior, so the result is not an isolated crossing. For the same reason, the result survives the multiplicity of testing 15 values of $k$. A fixed-sequence procedure, testing $k = 15, 14, \ldots$ in order and stopping at the first failure, controls the familywise error rate without adjustment and yields the same $k^*$. The sensitivity margins gave $k^* = \KStarMarginPtOne$ at $\delta = 0.01$ and $k^* = \KStarMarginPtThree$ at $\delta = 0.03$. Removing four electrodes (\KStarReductionPct\% of the array) is therefore within the prespecified tolerance. Larger reductions are not: halving the array to $k = 8$ cost $\DeltaKEight$ (95\% CI $\DeltaKEightLCB$ to $\DeltaKEightUCB$), and $k = 4$ cost $\DeltaKFour$. All paired differences at $k < 16$ were negative. The claim is non-inferiority within a stated engineering tolerance, not equivalence.

\begin{table*}[t]
\centering
\small
\caption{Sparse sensor curve and non-inferiority (cross-fitted ranking, $\lambda = 0$; 43 participants; every $k$ retrained from scratch on a fixed 200k-window outer-training subsample). $\Delta_k$ = paired difference to the $k = 16$ model of the same curve. Non-inferior when the lower 95\% bound exceeds $-\delta$; $\delta = 0.02$ is primary.}
\label{tab:sparse_ni}
\begin{tabular}{rcccccc}
\toprule
$k$ & Macro-F1 (mean $\pm$ SD) & $\Delta_k$ & 95\% CI & $\delta$=0.01 & $\delta$=0.02 & $\delta$=0.03 \\
\midrule
16 & 0.541 $\pm$ 0.084 & reference &  &  &  &  \\
15 & 0.538 $\pm$ 0.083 & $-$0.0036 & $-$0.0056 to $-$0.0017 & yes & yes & yes \\
14 & 0.535 $\pm$ 0.083 & $-$0.0060 & $-$0.0081 to $-$0.0039 & yes & yes & yes \\
13 & 0.530 $\pm$ 0.083 & $-$0.0109 & $-$0.0132 to $-$0.0085 & no & yes & yes \\
12 & 0.528 $\pm$ 0.084 & $-$0.0132 & $-$0.0160 to $-$0.0105 & no & yes & yes \\
11 & 0.521 $\pm$ 0.084 & $-$0.0202 & $-$0.0245 to $-$0.0160 & no & no & yes \\
10 & 0.515 $\pm$ 0.083 & $-$0.0258 & $-$0.0311 to $-$0.0205 & no & no & no \\
9 & 0.510 $\pm$ 0.086 & $-$0.0316 & $-$0.0385 to $-$0.0245 & no & no & no \\
8 & 0.505 $\pm$ 0.085 & $-$0.0366 & $-$0.0440 to $-$0.0289 & no & no & no \\
7 & 0.500 $\pm$ 0.085 & $-$0.0412 & $-$0.0488 to $-$0.0331 & no & no & no \\
6 & 0.480 $\pm$ 0.083 & $-$0.0609 & $-$0.0711 to $-$0.0507 & no & no & no \\
5 & 0.428 $\pm$ 0.076 & $-$0.1132 & $-$0.1276 to $-$0.0993 & no & no & no \\
4 & 0.397 $\pm$ 0.071 & $-$0.1444 & $-$0.1590 to $-$0.1306 & no & no & no \\
3 & 0.340 $\pm$ 0.057 & $-$0.2014 & $-$0.2211 to $-$0.1820 & no & no & no \\
2 & 0.299 $\pm$ 0.057 & $-$0.2420 & $-$0.2603 to $-$0.2230 & no & no & no \\
1 & 0.181 $\pm$ 0.040 & $-$0.3602 & $-$0.3831 to $-$0.3368 & no & no & no \\
\bottomrule
\end{tabular}
\end{table*}


\begin{figure*}[t]
\centering
\includegraphics[width=\textwidth]{figures/fig4_sparse_curve.pdf}
\caption{Sparse sensor curve (43 participants). (a) Mean Macro-F1 of models retrained on the top-$k$ electrodes of each fold's cross-fitted ranking, for the four values of $\lambda$ (curves overlap). Band = 95\% cluster-bootstrap CI for $\lambda = 0$. (b) Paired difference to the $k = 16$ model with 95\% CI. Dashed lines mark the prespecified margin ($\delta = 0.02$) and the two sensitivity margins.}
\label{fig:curve}
\end{figure*}

\subsection{H2: attribution-ranked subsets beat random subsets; other rankings match them}
Under the shared comparison model (Table~\ref{tab:baselines}, Fig.~\ref{fig:baselines}), SR-CFA subsets outperformed the mean of random subsets at every tested $k$ (Holm-corrected, 8 tests). The advantage was largest at intermediate sizes: $\HTwoKSixDiff$ Macro-F1 at $k = 6$ (95\% CI $\HTwoKSixLCB$ to $\HTwoKSixUCB$, $p_{\mathrm{Holm}} \HTwoKSixP$, $r_{rb} = \HTwoKSixR$) and $\HTwoKEightDiff$ at $k = 8$. It was small at $k = 4$ ($\HTwoKFourDiff$) and at $k = 12$ ($\HTwoKTwelveDiff$, $p_{\mathrm{Holm}} \HTwoKTwelveP$). Within each participant's own random subsets, the SR-CFA subset had a median percentile of \PctKSix\ at $k = 6$, \PctKEight\ at $k = 8$, \PctKFour\ at $k = 4$ and \PctKTwelve\ at $k = 12$. It reached the 95th percentile for only \PctKAboveNinetyFiveSix\ of 43 participants at $k = 6$ and \PctKAboveNinetyFiveTwelve\ at $k = 12$. Attribution therefore reliably beats the average random subset, but it rarely finds the best subset available for a given person.

In exploratory contrasts, attribution was not better than simpler alternatives. It did not differ from permutation importance at $k = 4$, 8 or 12 (e.g.\ $k = 8$: $\VSPermKEightDiff$, $p \VSPermKEightP$), and it exceeded it only at $k = 6$ ($\VSPermKSixDiff$, $p \VSPermKSixP$). The geometric baseline, which spreads electrodes evenly around both rings without using any data, matched SR-CFA at $k = 4$ ($\VSGeomKFourDiff$, $p \VSGeomKFourP$), $k = 8$ ($\VSGeomKEightDiff$, $p \VSGeomKEightP$) and $k = 12$ ($\VSGeomKTwelveDiff$, $p \VSGeomKTwelveP$). It was worse only at $k = 6$ ($\VSGeomKSixDiff$, $p \VSGeomKSixP$). Mutual information was the weakest ranking at every $k \le 8$.

\begin{table*}[t]
\centering
\small
\caption{Channel-selection methods at matched $k$ (43 participants). All methods use one comparison model (XGBoost, 50 trees, 150k-window outer-training subsample), so values are comparable within the table but lower than Table~\ref{tab:sparse_ni}. H2 $p$ values Holm-corrected over 8 tests; $^\dagger$exploratory, uncorrected.}
\label{tab:baselines}
\begin{tabular}{lcccc}
\toprule
Method & $k$ = 4 & $k$ = 6 & $k$ = 8 & $k$ = 12 \\
\midrule
SR-CFA ($\lambda$=0.5) & 0.383 $\pm$ 0.071 & 0.453 $\pm$ 0.080 & 0.469 $\pm$ 0.081 & 0.485 $\pm$ 0.079 \\
Cross-fitted SHAP ($\lambda$=0) & 0.382 $\pm$ 0.071 & 0.453 $\pm$ 0.080 & 0.468 $\pm$ 0.081 & 0.485 $\pm$ 0.079 \\
Permutation importance & 0.385 $\pm$ 0.070 & 0.440 $\pm$ 0.080 & 0.466 $\pm$ 0.080 & 0.485 $\pm$ 0.079 \\
Geometric farthest-point & 0.383 $\pm$ 0.069 & 0.419 $\pm$ 0.068 & 0.465 $\pm$ 0.072 & 0.478 $\pm$ 0.076 \\
Mutual information & 0.280 $\pm$ 0.049 & 0.360 $\pm$ 0.066 & 0.396 $\pm$ 0.068 & 0.480 $\pm$ 0.079 \\
Random subsets (mean) & 0.366 $\pm$ 0.059 & 0.416 $\pm$ 0.068 & 0.445 $\pm$ 0.071 & 0.478 $\pm$ 0.077 \\
\midrule SR-CFA $-$ random (H2) & +0.0172, $p$=0.007 & +0.0370, $p$=$<$0.001 & +0.0243, $p$=$<$0.001 & +0.0074, $p$=$<$0.001 \\
SR-CFA $-$ perm$^\dagger$ & $-$0.0018, $p$=0.678 & +0.0126, $p$=0.009 & +0.0033, $p$=0.069 & +0.0001, $p$=0.995 \\
SR-CFA $-$ geom$^\dagger$ & $-$0.0005, $p$=0.986 & +0.0341, $p$=$<$0.001 & +0.0047, $p$=0.267 & +0.0070, $p$=0.058 \\
SR-CFA percentile in random draws (median) & 65 (100 draws) & 89 (100 draws) & 82 (100 draws) & 66 (500 draws) \\
\bottomrule
\end{tabular}
\end{table*}


\begin{figure*}[t]
\centering
\includegraphics[width=\textwidth]{figures/fig5_selection_baselines.pdf}
\caption{Selection methods at matched $k$ (43 participants; comparison model shared by all methods). (a) Mean Macro-F1 with 95\% cluster-bootstrap CI. (b) Percentile of the SR-CFA subset within each participant's own random subsets (one point per participant; bar = median; 100 draws, 500 at $k = 12$).}
\label{fig:baselines}
\end{figure*}

\subsection{H3: the stability penalty had no measurable effect}
The instability term was small relative to attribution magnitude: the median $U_c/A_c$ was \UOverAMedian\% (maximum \UOverAMax\%). Penalising it therefore rarely changed a ranking. Against $\lambda = 0$, the top-12 set changed for \TopTwelveDiffersLambdaPtTwoFive, \TopTwelveDiffersLambdaPtFive\ and \TopTwelveDiffersLambdaOne\ of 43 folds at $\lambda = 0.25$, 0.5 and 1. Across all four $k$ tested, at least \HThreeMinTied\ of 43 participants received the identical subset under both rankings. None of the \HThreeNTests\ contrasts was significant after Holm correction (largest mean difference \HThreeMaxAbsDiff; smallest $p_{\mathrm{Holm}} \HThreeMinP$; Supplementary Table~\ref{tab:lambda}). The four curves in Fig.~\ref{fig:curve}a are indistinguishable (maximum spread between $\lambda$ values \LambdaCurveMaxSpread), and $k^*$ was identical under every $\lambda$ and margin. H3 is not supported. In these data, the cross-fitted magnitude ranking is already stable enough that a stability penalty has nothing to correct.

\subsection{H4: the attribution is faithful}
When all three deletion arms were trained under one protocol, the expected ordering held at every retained size (Table~\ref{tab:faithfulness}, Fig.~\ref{fig:faithfulness}). Deleting the highest-ranked electrodes hurt more than deleting random ones, and random deletion hurt more than deleting the lowest-ranked. \FaithNSig\ of the \FaithNTests\ Holm-corrected contrasts were significant. At six retained electrodes, top-removed scored \FaithTopKSix, random \FaithRandomKSix\ and bottom-removed \FaithBottomKSix\ (top $-$ random $\FaithTRKSixDiff$, $p_{\mathrm{Holm}} \FaithTRKSixP$; random $-$ bottom $\FaithRBKSixDiff$, $p_{\mathrm{Holm}} \FaithRBKSixP$). The only non-significant contrast was top- versus bottom-removal at 14 retained electrodes ($\FaithTBKFourteenDiff$, $p_{\mathrm{Holm}} \FaithTBKFourteenP$), where only two electrodes are removed. The gap between arms was largest at intermediate sizes and small at the extremes. With few electrodes left all arms perform poorly; with many left, the array is redundant enough that any two electrodes can be spared.

\begin{figure}[t]
\centering
\includegraphics[width=\columnwidth]{figures/fig7_faithfulness.pdf}
\caption{Deletion-with-retraining test of attribution faithfulness (43 participants; mean and 95\% cluster-bootstrap CI). All arms use the same comparison model and training subsample. Random removal is available where random subsets were sampled ($n$ = 4, 6, 8, 12).}
\label{fig:faithfulness}
\end{figure}

\begin{table*}[t]
\centering
\small
\caption{Deletion (remove-and-retrain) faithfulness test (H4; 43 participants). Every arm uses the same comparison model and training subsample; random removal averages 100 (500 at $k$ = 12) draws per participant and exists only where those draws were sampled. Differences are paired means with Holm-corrected $p$ in parentheses.}
\label{tab:faithfulness}
\begin{tabular}{rcccccc}
\toprule
Kept & Top removed & Random & Bottom removed & Top$-$Bottom & Top$-$Random & Random$-$Bottom \\
\midrule
2 & 0.263 & -- & 0.292 & $-$0.0297 ($<$0.001) & -- & -- \\
4 & 0.351 & 0.366 & 0.382 & $-$0.0314 ($<$0.001) & $-$0.0151 (0.010) & $-$0.0163 (0.013) \\
6 & 0.376 & 0.416 & 0.454 & $-$0.0775 ($<$0.001) & $-$0.0398 ($<$0.001) & $-$0.0377 ($<$0.001) \\
8 & 0.408 & 0.445 & 0.469 & $-$0.0618 ($<$0.001) & $-$0.0377 ($<$0.001) & $-$0.0242 ($<$0.001) \\
10 & 0.445 & -- & 0.475 & $-$0.0297 ($<$0.001) & -- & -- \\
12 & 0.455 & 0.478 & 0.485 & $-$0.0302 ($<$0.001) & $-$0.0228 ($<$0.001) & $-$0.0074 ($<$0.001) \\
14 & 0.487 & -- & 0.491 & $-$0.0038 (0.150) & -- & -- \\
\bottomrule
\end{tabular}
\end{table*}


\subsection{Stability of the ranking across people}
Rankings computed from individual participants' attributions agreed closely with each other (Table~\ref{tab:stability}): mean Spearman $\rho$ = \StabSubjRho\ $\pm$ \StabSubjRhoSD\ over \StabSubjPairs\ pairs, and Jaccard overlap of the top six = \StabSubjJacSix. Rankings from disjoint inner folds agreed at $\rho$ = \StabInnerRho. The outer-fold agreement ($\rho$ = \StabOuterRho) is an upper bound only, because those training sets share about 98\% of their participants. Across folds, the highest- and lowest-ranked electrodes kept their positions, and the rank ranges were widest in the middle of the ordering (Fig.~\ref{fig:attribution}b).

\begin{table*}[t]
\centering
\small
\caption{Agreement of channel rankings ($\lambda = 0$). Values are means over pairs.}
\label{tab:stability}
\begin{tabular}{lcccccc}
\toprule
Comparison & Pairs & Spearman $\rho$ & Kendall $\tau$ & Jaccard@4 & Jaccard@6 & Jaccard@8 \\
\midrule
Per-subject attribution (between people) & 903 & 0.927 $\pm$ 0.038 & 0.801 & 0.752 & 0.905 & 0.803 \\
Inner folds (disjoint subject groups) & 430 & 0.948 $\pm$ 0.021 & 0.840 & 0.816 & 0.992 & 0.838 \\
Outer folds (upper bound$^*$) & 903 & 0.987 $\pm$ 0.008 & 0.943 & 0.888 & 1.000 & 0.923 \\
\bottomrule
\end{tabular}
\par\smallskip\parbox{0.95\linewidth}{\footnotesize $^*$Outer training sets share about 98\% of their participants, so their agreement is an upper bound, not an estimate of between-subject stability.}
\end{table*}


\begin{figure*}[t]
\centering
\includegraphics[width=\textwidth]{figures/fig6_attribution.pdf}
\caption{Cross-fitted channel attribution ($\lambda = 0$, 43 outer folds). (a) Mean share of attribution per electrode, drawn on the $2 \times 8$ montage (ring~2 electrodes are aligned with, and 2~cm from, ring~1 electrodes at the same position). (b) Rank of each electrode across folds (median and range).}
\label{fig:attribution}
\end{figure*}

\subsection{Exploratory analyses}

\begin{figure*}[t]
\centering
\includegraphics[width=\textwidth]{figures/fig8_generalisation.pdf}
\caption{Exploratory analyses. (a) Cross-session: each point is one held-out participant (bar = median); all conditions are also subject-independent. (b) NinaPro DB5 sparse curve (10 subjects; mean and 95\% cluster-bootstrap CI; dotted line = chance). (c) GRABMyo bipolar-pair sparse curve (43 participants; $\lambda = 0$).}
\label{fig:generalisation}
\end{figure*}

\paragraph{Cross-session}
Training on session~1 and testing on sessions 1, 2 and 3 gave \CSOneSOne, \CSOneSTwo\ and \CSOneSThree\ (Table~\ref{tab:cross_session}, Fig.~\ref{fig:generalisation}a). Performance on later sessions was not lower than on the training session. The session-2 score was higher ($\CSSTwoVsSOneDiff$, 95\% CI $\CSSTwoVsSOneLCB$ to $\CSSTwoVsSOneUCB$, uncorrected $p \CSSTwoVsSOneP$), and session 3 did not differ ($\CSSThreeVsSOneDiff$, $p \CSSThreeVsSOneP$). Because every condition is also subject-independent, this does not show that sessions are interchangeable. It shows that, once the cross-user gap is present, no additional cross-day loss was detectable. Training on sessions 1 and 2 improved session-3 performance over training on session 1 alone ($\CSTwoVsOneDiff$, 95\% CI $\CSTwoVsOneLCB$ to $\CSTwoVsOneUCB$, $p \CSTwoVsOneP$). This design cannot separate day diversity from the doubled amount of training data.

\paragraph{Bipolar montage}
With the eight derived bipolar pairs as the selection unit, full-array LOSO Macro-F1 was \BipFull\ $\pm$ \BipFullSD, lower than the monopolar array (Fig.~\ref{fig:generalisation}c, Supplementary Table~\ref{tab:replication}). No pair could be removed within the primary margin. With seven pairs, $\bar\Delta_7 = \BipDeltaKSeven$, and its lower bound ($\BipDeltaKSevenLCB$) fell just outside $-0.02$. It did so under every one of 200 alternative bootstrap seeds (range $\BipSeedLCBMinSeven$ to $\BipSeedLCBMaxSeven$). $k^*$ was therefore \BipKStarMarginPtTwo\ at $\delta = 0.02$, \BipKStarMarginPtOne\ at 0.01, and \BipKStarMarginPtThree\ pairs at 0.03, identically for all $\lambda$. The ranking was again stable: pairs \BipTopThreeSet\ formed the top three in \BipTopThreeFolds\ of 43 folds. In deployable sensor units, then, the array showed no redundancy within the prespecified tolerance.

\paragraph{NinaPro DB5}
On DB5 the full 16-channel LOSO Macro-F1 was \DBFiveFull\ $\pm$ \DBFiveFullSD\ (12 movements; chance about 0.083), far below GRABMyo and with wide between-subject spread (\DBFiveFullMin--\DBFiveFullMax) (Fig.~\ref{fig:generalisation}b, Supplementary Table~\ref{tab:replication}). Under the frozen rule, $k^* = \DBFiveKStarMarginPtTwo$ at $\delta = 0.02$ ($\bar\Delta_{12} = \DBFiveDeltaKTwelve$, lower bound $\DBFiveDeltaKTwelveLCB$). This result is fragile. The lower bound lies on the margin, and $k = 12$ was non-inferior under \DBFiveSeedFrac\% of \DBFiveSeedN\ alternative bootstrap seeds (lower bound range $\DBFiveSeedLCBMin$ to $\DBFiveSeedLCBMax$). Moreover, $k = 13$, 14 and 15 all failed (lower bounds $\DBFiveDeltaKThirteenLCB$, $\DBFiveDeltaKFourteenLCB$, $\DBFiveDeltaKFifteenLCB$). A fixed-sequence procedure starting at $k = 15$ would stop immediately and report no non-inferior subset. With ten subjects, DB5 does not support a claim that a quarter of the channels can be removed. The procedure itself transferred without modification, and the ranked curve had the same monotone shape.

\paragraph{Fixed layout (post hoc)}
One 12-electrode layout (removing electrodes \FixedDropped) was derived on \FixedDev\ participants and applied unchanged to \FixedTest\ others. It scored $\FixedDiff$ against the 16-electrode model trained on the same participants (95\% CI $\FixedLCB$ to $\FixedUCB$). This lies within the margin, but it is one post-hoc split and is reported only as a consistency check on the LOSO result.

\paragraph{Feature redundancy}
Removing IEMG, an exact multiple of MAV, changed Macro-F1 by $\RedundDiff$ (95\% CI $\RedundLCB$ to $\RedundUCB$; 43 folds, comparison model). Channel rankings were almost unchanged: the median per-fold Spearman $\rho$ between rankings with and without IEMG was \RedundRho, and the top-12 set was identical in \RedundTopTwelveSame\ of 43 folds. The duplicated feature therefore did not distort the channel ranking.

\section{Discussion}
\label{sec:discussion}

\paragraph{What the prespecified tests show}
Under a margin fixed before the curve was computed, four of the 16 GRABMyo electrodes could be removed while remaining non-inferior to the full array. Five could not. Every $k$ from 12 to 15 met the margin, so the result does not rest on a single lucky crossing. The margin is an engineering tolerance, not a clinical threshold. Every reduction cost some Macro-F1, and halving the array cost about \DeltaKEight. The practical reading is modest: in a subject-independent setting, this 16-electrode array contains only limited redundancy for 16-gesture recognition. User-dependent studies that retain far fewer channels answer a different question~\cite{phinyomark2018emg,qu2021reduce}, and their results should not be taken as evidence that a population-level layout can be equally sparse.

Two of the three hypotheses about attribution were supported. Attribution-ranked subsets beat the average random subset at every size tested (H2), and deletion with retraining showed the expected ordering at every size (H4). The ranking is therefore not decorative: it reflects what the retrained models depend on. The stability penalty, however, did nothing (H3). With cross-fitting over roughly 34 participants per inner-training set, inner-fold attributions already agreed so closely that the dispersion term was a few percent of the magnitude term. It could not reorder channels in a way that mattered. We report SR-CFA as a tested and unsupported variant. The supported method is plain cross-fitted TreeSHAP magnitude, and a related robustness-weighted score has been proposed independently~\cite{xiang2026roshap}.

\paragraph{Attribution is not better than simple alternatives}
The exploratory comparisons temper H2. Beating random subsets is a low bar. A data-free geometric rule that spreads electrodes around both rings matched the attribution ranking at 4, 8 and 12 electrodes, and permutation importance matched it at three of four sizes. Attribution had a clear advantage only at $k = 6$. Within each participant's own random subsets, the attribution subset was typically above the median but rarely near the top. This fits a picture in which even circumferential coverage explains most of what any ranking can achieve, and the choice among well-spread layouts matters little. For a device designer, the evidence favours a well-spread layout of about 12 electrodes over trusting any particular ranking to identify an unusually good sparse subset.

\paragraph{Interpretation of the ranked electrodes}
The four highest-ranked electrodes (\TopFourSet) were the same in \TopFourFolds\ of 43 folds and in single-participant rankings (mean Spearman $\rho$ = \StabSubjRho). We deliberately do not map them to named muscles. The dataset documents ring positions, not the anatomical placement of each circumferential index. Surface electrodes also record compound activity from several muscles at each site~\cite{farina2004extraction}. The robust statement is positional. High attribution concentrated at a few circumferential positions (Fig.~\ref{fig:attribution}a). The top and bottom of the ranking were stable, and the ordering varied mainly among mid-ranked electrodes. This fits the finding that many different well-spread subsets perform similarly.

\paragraph{Cross-session and replication}
Cross-session results were exploratory and should be read narrowly. Held-out participants were not recognised worse on later days than on the training day. This shows only that no extra cross-day loss was detectable on top of the cross-user gap, not that electrode shift is harmless~\cite{hargrove2006effect,kaufmann2010fluctuating,vidovic2016improving}. The gain from training on two sessions is confounded with doubling the training data. On DB5 the whole procedure ran without modification, but the replication was underpowered. Ten subjects and a low baseline left the $k = 12$ decision on the margin, with neighbouring $k$ failing. It is best read as a non-contradiction rather than a replication of the reduction.

\paragraph{Limitations}
The primary unit is a monopolar electrode. A conventional device measures bipolar pairs, and removing four electrodes does not remove four sensors of such a device. The exploratory bipolar re-analysis points the other way: with pairs as the unit, not even one of eight could be removed within the primary margin, although the shortfall at seven pairs was small. The 25\% reduction should therefore be read as a statement about monopolar electrodes on this array, not as evidence that a bipolar armband can shed a quarter of its sensors. The bipolar reference is itself a derived montage, and physical bipolar electrodes may behave differently. All participants were able-bodied, and all analysis was offline. The SVM was trained on a subsample and is a lower bound. The H2 and H4 comparisons used a lighter comparison model than the headline curve. The count-feature threshold was a single scalar set from pooled data. $\lambda$ was not selected, and the detailed H2--H4 test families were specified after the underlying results were computed (Section~\ref{sec:deviations}). Cross-model attribution, window-length sensitivity and the 17-class task were listed as exploratory in the plan and were not run.

\section{Conclusion}
\label{sec:conclusion}

We evaluated explainability-guided electrode reduction for subject-independent sEMG gesture recognition using cross-fitted attribution and decision rules fixed in advance. On 43 GRABMyo participants, a 12-of-16 electrode subset was non-inferior to the full array within a 0.02 Macro-F1 engineering tolerance; smaller subsets were not. Cross-fitted TreeSHAP rankings were stable across people. They beat random subsets and were faithful under deletion with retraining. A stability penalty added nothing, and a data-free, evenly spread layout performed about as well as the attribution ranking at most sizes. The results support moderate, not aggressive, electrode reduction for population-level devices. In an exploratory re-analysis with bipolar pairs as the unit, no pair could be removed within the same tolerance. They also suggest that the benefit of attribution lies in confirming which positions matter, rather than in finding sparse layouts that simple coverage rules would miss.


\section*{Declaration of Competing Interest}
The authors declare that they have no known competing financial interests or personal relationships that could have appeared to influence the work reported in this paper.

\section*{Data and Code Availability}
GRABMyo v1.1.0 is available from PhysioNet (\url{https://physionet.org/content/grabmyo/1.1.0/}) and NinaPro DB5 from the NinaPro repository. The frozen analysis plan, decision log, analysis code and every result file cited in this paper are included in the accompanying repository. All numbers in the text and tables are generated from those files.

\appendix
\section{Supplementary tables}
\begin{table}[t]
\centering
\small
\caption{Stability regularisation versus attribution magnitude alone (H3): $Q(\lambda)$ ranking minus $Q(0)$ ranking on the sparse curve. Tied = participants for whom both rankings select the identical subset. Holm over 12 tests.}
\label{tab:lambda}
\begin{tabular}{rrcccc}
\toprule
$\lambda$ & $k$ & Mean diff & 95\% CI & Tied & $p_{\mathrm{Holm}}$ \\
\midrule
0.25 & 4 & $-$0.0002 & $-$0.0005 to +0.0000 & 41/43 & 1.000 \\
0.25 & 6 & +0.0000 & +0.0000 to +0.0000 & 43/43 & 1.000 \\
0.25 & 8 & +0.0001 & +0.0000 to +0.0002 & 42/43 & 1.000 \\
0.25 & 12 & +0.0000 & +0.0000 to +0.0000 & 43/43 & 1.000 \\
0.5 & 4 & +0.0005 & $-$0.0004 to +0.0019 & 39/43 & 1.000 \\
0.5 & 6 & $-$0.0001 & $-$0.0002 to +0.0000 & 42/43 & 1.000 \\
0.5 & 8 & +0.0008 & +0.0000 to +0.0020 & 40/43 & 1.000 \\
0.5 & 12 & $-$0.0000 & $-$0.0001 to +0.0000 & 42/43 & 1.000 \\
1 & 4 & +0.0010 & $-$0.0002 to +0.0028 & 37/43 & 1.000 \\
1 & 6 & $-$0.0001 & $-$0.0002 to +0.0000 & 42/43 & 1.000 \\
1 & 8 & +0.0015 & +0.0004 to +0.0029 & 36/43 & 0.216 \\
1 & 12 & $-$0.0005 & $-$0.0014 to +0.0000 & 41/43 & 1.000 \\
\bottomrule
\end{tabular}
\end{table}

\begin{table*}[t]
\centering
\small
\caption{Cross-session evaluation (exploratory). All conditions are also subject-independent: the test participant is never in training.}
\label{tab:cross_session}
\begin{tabular}{lcccc}
\toprule
Design & Train sessions & Test session & Train windows & Macro-F1 \\
\midrule
CS1 & 1 & 1 & 183{,}456 & 0.492 $\pm$ 0.101 \\
CS1 & 1 & 2 & 183{,}456 & 0.528 $\pm$ 0.088 \\
CS1 & 1 & 3 & 183{,}456 & 0.510 $\pm$ 0.090 \\
CS2 & 1+2 & 3 & 366{,}912 & 0.530 $\pm$ 0.095 \\
\bottomrule
\end{tabular}
\par\smallskip\parbox{0.95\linewidth}{\footnotesize Paired contrasts (exploratory, uncorrected): CS1: test S2 $-$ test S1: +0.0355 (95\% CI +0.0091 to +0.0621), $p$ = 0.024; CS1: test S3 $-$ test S1: +0.0172 (95\% CI $-$0.0038 to +0.0391), $p$ = 0.272; S3 test: CS2 (train S1+S2) $-$ CS1 (train S1): +0.0210 (95\% CI +0.0146 to +0.0274), $p$ $<$ 0.001.}
\end{table*}

\begin{table*}[t]
\centering
\small
\caption{Replications of the sparse curve (cross-fitted ranking, $\lambda = 0$; margin $\delta = 0.02$). Left: NinaPro DB5, 10 subjects, reference $k = 16$; seed column = share of 200 alternative bootstrap seeds under which the $k$ is non-inferior. Right: GRABMyo bipolar pairs $F_j - F_{j+8}$, 43 participants, reference $k = 8$.}
\label{tab:replication}
\begin{tabular}{rccccccc}
\toprule
$k$ & DB5 F1 & $\Delta_k$ & LCB & seeds NI & Bipolar F1 & $\Delta_k$ & LCB \\
\midrule
15 & 0.270 & $-$0.0108 & $-$0.0238 & 0\% &  &  &  \\
14 & 0.269 & $-$0.0119 & $-$0.0274 & 0\% &  &  &  \\
13 & 0.266 & $-$0.0143 & $-$0.0312 & 0\% &  &  &  \\
12 & 0.268 & $-$0.0126 & $-$0.0200 & 96\% &  &  &  \\
11 & 0.260 & $-$0.0211 & $-$0.0341 & 0\% &  &  &  \\
10 & 0.265 & $-$0.0160 & $-$0.0314 & 0\% &  &  &  \\
9 & 0.260 & $-$0.0202 & $-$0.0369 & 0\% &  &  &  \\
8 & 0.252 & $-$0.0288 & $-$0.0554 & 0\% &  &  &  \\
7 & 0.232 & $-$0.0486 & $-$0.0753 & 0\% & 0.463 & $-$0.0155 & $-$0.0203 \\
6 & 0.231 & $-$0.0498 & $-$0.0818 & 0\% & 0.441 & $-$0.0376 & $-$0.0456 \\
5 & 0.200 & $-$0.0803 & $-$0.1184 & 0\% & 0.405 & $-$0.0737 & $-$0.0841 \\
4 & 0.180 & $-$0.1006 & $-$0.1503 & 0\% & 0.354 & $-$0.1245 & $-$0.1395 \\
3 & 0.162 & $-$0.1185 & $-$0.1741 & 0\% & 0.309 & $-$0.1703 & $-$0.1857 \\
2 & 0.138 & $-$0.1430 & $-$0.1953 & 0\% & 0.239 & $-$0.2394 & $-$0.2608 \\
1 & 0.105 & $-$0.1754 & $-$0.2295 & 0\% & 0.174 & $-$0.3047 & $-$0.3295 \\
\bottomrule
\end{tabular}
\end{table*}


\bibliographystyle{elsarticle-num}
\bibliography{references}

\end{document}
